\documentclass[11pt,a4paper]{article}

\usepackage{kotex}
\usepackage{fontspec}
\setmainfont{Times New Roman}
\setmainhangulfont{AppleMyungjo}
\setsanshangulfont{Apple SD Gothic Neo}

\usepackage[a4paper,top=14mm,bottom=16mm,left=20mm,right=20mm]{geometry}
\usepackage{fancyhdr}
\setlength{\headheight}{24pt}
\usepackage{enumitem}
\usepackage{mdframed}
\usepackage{array}
\usepackage{booktabs}

\setlength{\parindent}{1em}
\linespread{1.04}

\pagestyle{fancy}
\fancyhf{}
\renewcommand{\headrulewidth}{0.6pt}
\fancyhead[L]{\sffamily\bfseries 2026 고1 영어 변형 — 지문 27}
\fancyhead[R]{\sffamily [Junior Coding Camp 안내문]}
\renewcommand{\footrulewidth}{0pt}
\fancyfoot[C]{\framebox[16mm][c]{\thepage}}

\newcommand{\qno}[1]{\par\vspace{5pt}\noindent{\sffamily\bfseries #1.}\ }
\newcommand{\instr}[1]{{\sffamily #1}\par\vspace{2pt}}
\newlist{choices}{enumerate}{1}
\setlist[choices]{label=,leftmargin=1.5em,itemsep=1pt,topsep=2pt,parsep=0pt}
\newmdenv[linewidth=0.6pt,roundcorner=3pt,innertopmargin=7pt,innerbottommargin=7pt,
          innerleftmargin=9pt,innerrightmargin=9pt,skipabove=5pt,skipbelow=5pt]{pbox}
\newcommand{\nemo}[1]{\fbox{\,#1\,}}

\begin{document}

\begin{center}
{\fontsize{15}{17}\selectfont\sffamily\bfseries [지문 27] 다음 글을 읽고 물음에 답하시오.}
\end{center}
\vspace{4pt}

\begin{pbox}
\begin{center}
{\large\bfseries Junior Coding Camp}\\[4pt]
{\itshape Explore the world of coding through easy, hands-on activities!}
\end{center}

\vspace{6pt}
\textbf{Date \& Site}
\begin{itemize}[leftmargin=1.8em,itemsep=1pt,topsep=2pt,parsep=0pt]
  \item[\textbullet] Wednesday, July 29th, 2026
  \item[\textbullet] Arendale Tech Center
\end{itemize}

\vspace{4pt}
\textbf{Ages \& Level:} 10--13 years old, for beginners

\vspace{4pt}
\textbf{Schedule}

\vspace{2pt}
\begin{tabular}{@{\hspace{1em}}ll}
\toprule
\textbf{Time} & \textbf{Activity} \\
\midrule
9 a.m. -- 12 p.m. & Coding Basics \\
1 p.m. -- 4 p.m.  & Game Design \\
\bottomrule
\end{tabular}

\vspace{4pt}
\textbf{Participation Fee:} \$100 per person (lunch provided)

\vspace{4pt}
\textbf{Registration}
\begin{itemize}[leftmargin=1.8em,itemsep=1pt,topsep=2pt,parsep=0pt]
  \item[\textbullet] Limited to 50 students
  \item[\textbullet] On a first-come, first-served basis
  \item[\textbullet] Register on site.
\end{itemize}
\end{pbox}

\vspace{6pt}

%% Q1 — 내용 불일치 (2점)
\qno{1}\instr{윗글의 내용과 일치하지 않는 것은?}
\begin{choices}
  \item ① The event spans a single day at a technology facility.
  \item ② Participation is restricted to those with no prior coding experience.
  \item ③ Those who register early are given priority for spots.
  \item ④ Participants are required to bring their own meals as no food is provided.
  \item ⑤ The morning session focuses on foundational coding skills.
\end{choices}

\vspace{10pt}

%% Q2 — 서술형 단답 (2점)
\qno{2}\instr{[서술형] 다음 물음에 답하시오.}

\noindent (1)\ 캠프에 참가하기 위한 연령 조건을 우리말로 쓰시오.

\vspace{4pt}
\noindent $\Rightarrow$ \rule{10cm}{0.4pt}

\vspace{10pt}

\noindent (2)\ 등록 방법(where to register)을 본문에서 찾아 영어로 쓰시오.

\vspace{4pt}
\noindent $\Rightarrow$ \rule{10cm}{0.4pt}

% ===================== Q3 =====================
\qno{3}\instr{오후 일정에 해당하는 활동으로 가장 적절한 것은?}
\begin{choices}
  \item ① Coding Basics
  \item ② Game Design
  \item ③ Lunch preparation
  \item ④ Online registration
  \item ⑤ Advanced robotics
\end{choices}

% ===================== Q4 =====================
\qno{4}\instr{윗글에서 추론할 수 있는 내용으로 가장 적절한 것은?}
\footnotesize
\begin{tabular}{@{}p{0.48\linewidth}p{0.48\linewidth}@{}}
① More than 50 students can join if they pay an additional fee. &
② Students with advanced coding experience are the main target group. \\
③ Applicants who register earlier have a better chance of being accepted. &
④ The camp lasts for two consecutive days. \\
⑤ Participants must find lunch outside the camp site. &
\end{tabular}
\normalsize

\end{document}
